% SPU-13: A Deterministic Rational-Field Processor for Exact Geometric Computation
% Central architecture paper — arXiv:cs.AR
%
% Build:  pdflatex spu13_central_paper.tex && bibtex spu13_central_paper && pdflatex spu13_central_paper.tex
%
% Status: Draft foundation paper. Keep silicon claims tied to evidence ledger.

\documentclass[10pt,twocolumn]{article}

\usepackage[hmargin=0.9in,vmargin=0.9in]{geometry}
\usepackage{cite}
\usepackage{amsmath,amssymb}
\usepackage{graphicx}
\usepackage{hyperref}
\usepackage{booktabs}
\usepackage{multirow}

\title{SPU-13: A Deterministic Rational-Field Processor \\
       for Exact Geometric Computation}

\author{John Curley \\
        Independent Researcher \\
        \texttt{github.com/johncurley/SPU}}

\begin{document}
\maketitle

\begin{abstract}
We present the SPU-13, an open-source FPGA coprocessor that performs all
geometric computation in the exact rational field $\mathbb{Q}(\sqrt{3})$,
eliminating floating-point drift entirely. The architecture combines a 13-axis
manifold engine, a Thimble-Pad\'e rational approximant pipeline over the
biquadratic extension field $A_{31}(\mathbb{M}_{31})$, and a
$\mathbb{Z}[\phi]/L_p$ phinary arithmetic co-processor for phase-coherent
operations. A compact SPI southbridge protocol from an RP2350-class
microcontroller provides instruction injection, configuration writes, result
readback, and USB telemetry. Current silicon evidence covers split subsystem
probes on Gowin GW5A-25A and SPI-visible integration spins on Xilinx
XC7A100T, including the live RPLU2 Pad\'e evaluator. Lattice ECP5 and larger
7-series Kintex devices are treated as portability and full-integration
targets rather than completed silicon proofs.
\end{abstract}

%==============================================================================
\section{Introduction}
\label{sec:introduction}
%==============================================================================

Any computation involving irrational constants -- $\pi$, $\sqrt{2}$, $\sqrt{5}$,
the golden ratio $\phi$ -- must either approximate them as finite-precision
floats or work in an algebraic extension where they are native. The former
accumulates error; the latter preserves exactness.\cite{wildberger2005}

The SPU-13 takes the latter approach entirely. Every register holds a pair
$(P,Q)$ representing the exact value $P+Q\sqrt{3}$. Arithmetic is closed in
this field: $\sqrt{3}\times\sqrt{3}=3$, a rational. No transcendental
functions, no division in hot paths, no floating-point units anywhere in the
core.

% - why Q(sqrt3) is the minimal closed field for IVM geometry
% - relationship to Buckminster Fuller's synergetics and Wildberger's rational trigonometry
% - the zero-drift invariant: any sequence of operations returns to the exact
%   starting bit pattern after a closed orbit

%==============================================================================
\section{Mathematical Foundation}
\label{sec:math}
%==============================================================================

\subsection{Rational Surd Fields and the IVM}
\label{sec:ivm}

The isotropic vector matrix (IVM) lattice has 60$^\circ$ angles between basis
vectors.\cite{fuller1975} The hexagonal cross-sections of this lattice involve
$\cos 60^\circ = 1/2$ and $\sin 60^\circ = \sqrt{3}/2$. Therefore, any
algebraic computation in the IVM encounters $\sqrt{3}$.

The field $\mathbb{Q}(\sqrt{3})$ is closed under multiplication:
\begin{equation}
(a+b\sqrt{3}) \times (c+d\sqrt{3}) = (ac+3bd) + (ad+bc)\sqrt{3}
\end{equation}

This means all quadrances (squared distances), all spreads, and all geometric
relationships in the IVM stay in $\mathbb{Q}(\sqrt{3})$ without ever escaping
into higher irrationals or transcendentals.\cite{wildberger2005}

\subsection{Pell Octave and the Rotor}
\label{sec:pell}

The Pell rotor $r = 2 + \sqrt{3}$ satisfies the Pell equation
$P^2 - 3Q^2 = 1$ for all integer powers. After 8 rotations the mantissa
resets and an octave counter increments, providing infinite rotor range in
16-bit registers.\cite{thomson2026}

\subsection{The Biquadratic Extension \texorpdfstring{$\mathbb{A}_{31}$}{A31}}
\label{sec:a31}

For the RPLU v2 pipeline, arithmetic is lifted to the biquadratic extension
field $\mathbb{A}_{31} = \mathbb{F}_{p^4}$ over the Mersenne prime
$p = 2^{31}-1$, with basis $[1, \sqrt{3}, \sqrt{5}, \sqrt{15}]$.
Multiplication in this field uses sixteen parallel $32\times32$ products
with Mersenne reduction via 72-bit chunk-split-and-add.

%==============================================================================
\section{System Architecture}
\label{sec:arch}
%==============================================================================

\subsection{Core Topology}
\label{sec:cores}

The repository is organized around a primary SPU-13 manifold core and smaller
satellite or sidecar cores used for targeted proofs:

\begin{table}[h]
\centering
\resizebox{\columnwidth}{!}{%
\begin{tabular}{lcc}
\toprule
Core & Axes & Role \\
\midrule
SPU-13 Cortex   & 13 & Manifold engine, RPLU pipeline, sidecar host \\
SPU-4 Sentinel  & 4  & Quadray satellite and reduced probe target \\
\bottomrule
\end{tabular}
}
\caption{Processor cores.}
\end{table}

\subsection{SPI Southbridge Protocol}
\label{sec:spi}

The current hardware control plane uses a unified command set over SPI mode 0.
Bench bring-up on Wukong Artix-7 uses conservative 25--100 kHz bit-banged
transactions on J11; the protocol is designed to move to faster SPI and then
PIO/DMA transport once the command semantics are frozen:

\begin{table}[h]
\centering
\resizebox{\columnwidth}{!}{%
\begin{tabular}{ccl}
\toprule
Command & Code & Function \\
\midrule
0xAC & Status    & Read FPGA status, telemetry \\
0xA0 & Manifold  & Burst read 13-axis manifold state \\
0xAE & QR Commit & Commit result to Quadray register file \\
0xB1 & Instruction & Write instruction (QLDI, QSUB, ROTC, etc.) \\
0xA5 & Config    & Write RPLU configuration record \\
\bottomrule
\end{tabular}
}
\caption{SPI southbridge command set.}
\end{table}

The protocol is transport-agnostic. Current hardware uses 4-wire SPI; a
parallel PIO transport using the same command set is in development for
higher throughput.

%==============================================================================
\section{Instruction Set}
\label{sec:isa}
%==============================================================================

The SPU-13 implements 26 opcodes across six functional groups:

\begin{table}[h]
\centering
\resizebox{\columnwidth}{!}{%
\begin{tabular}{ll}
\toprule
Group & Opcodes \\
\midrule
Quadray load/store  & QLDI, QSUB, QDISP \\
Rotors              & ROTC (6 corrected angles, period 2-6) \\
SOM/BMU             & SOM\_CLASSIFY, SOM\_TRAIN \\
RPLU config         & RPLU\_CFG, RPLU\_EVAL, RPLU\_CONSUME \\
Lucas phinary       & PSCALE, PCHIRAL, PMUL, PINV, PHSLK \\
Neuro sidecar       & NEURO\_CFG, NEURO\_START, NEURO\_SPIKE \\
\bottomrule
\end{tabular}
}
\caption{Opcode groups. See \cite{isa_reference} for full encoding.}
\end{table}

The arithmetic hot paths are branch-free and float-free. Multi-cycle sidecars
such as inversion use bounded finite-state machines with explicit busy/done
handshakes, so host-visible execution remains deterministic and auditable.

%==============================================================================
\section{RPLU v2 Thimble-Pad\'e Pipeline}
\label{sec:rplu}
%==============================================================================

The RPLU v2 pipeline is a four-stage rational approximant engine:

\begin{enumerate}
\item[$\Phi_1$] \textbf{SOM BMU} -- 7-node parallel Kohonen self-organizing
  map with weighted quadrance winner-take-all tree.
\item[$\Phi_2$] \textbf{BTU} -- Barycentric Transmutation Unit, routes the
  BMU saddle point into 4-lane BRAM for $\mathbb{A}_{31}$ evaluation.
\item[$\Phi_3$] \textbf{Thimble-Pad\'e} -- $[4/4]$ rational Pad\'e approximant
  evaluated via Horner's method, using the Conjugate Reduction Tower for
  $\mathbb{A}_{31}$ denominator inversion ($\sim$76 cycles, deterministic).
\item[$\Phi_4$] \textbf{Output latch} -- Final thimble contribution, quadray
  variety check, and QR commit.
\end{enumerate}

The pipeline uses a shared $\mathbb{A}_{31}$ multiplier with Mersenne
reduction. Detailed architecture and verification are in \cite{rplu2026}.

%==============================================================================
\section{Lucas Phinary Co-Processor}
\label{sec:lucas}
%==============================================================================

The Lucas MAC operates over $\mathbb{Z}[\phi]/L_{521}$, where
$\phi = (1+\sqrt{5})/2$ and $L_{13}=521$ is the 13th Lucas prime.

\begin{table}[h]
\centering
\begin{tabular}{lcc}
\toprule
Operation & Cycles & DSP \\
\midrule
PSCALE  ($\phi\cdot x$)        & 1 & 0 \\
PCHIRAL ($\overline{x}$)       & 1 & 0 \\
PMUL    ($x\cdot y$)           & 3 & 16 \\
PINV    ($x^{-1}$)             & $O(\log L_p)$ & 16 \\
PHSLK   (coherence predicate)  & 1 & 0 \\
\bottomrule
\end{tabular}
\caption{Lucas MAC opcodes.}
\end{table}

The PHSLK predicate checks rational phase coherence by cross-multiplication,
avoiding denominator inversion entirely. The software oracle performs a
million-step zero-drift marathon, Tang 25K silicon covers the PSCALE/PCHIRAL
fast paths, and Wukong Artix-7 silicon covers the SPI-visible PMUL/PINV
sidecar operations. Full details are in \cite{lucas2026}.

%==============================================================================
\section{Silicon Results}
\label{sec:results}
%==============================================================================

\subsection{Resource Utilization}
\label{sec:resources}

\begin{table}[h]
\centering
\resizebox{\columnwidth}{!}{%
\begin{tabular}{lrrrr}
\toprule
Spin & LUTs & FFs & DSP & BRAM \\
\midrule
Tang 25K: rplu2\_arith      & 9,211 & 7,926 & --- & --- \\
Tang 25K: rotc\_probe       & 13,352 & 1,036 & 0 & 0 \\
Tang 25K: lucas\_mac        & 696 & 216 & 0 & 0 \\
Artix-7: LUCAS              & 5,073 & --- & 120 & --- \\
Artix-7: SU3                & 21,092 & 9,488 & 64 & --- \\
Artix-7: RPLU2PADE          & 20,277 & 6,678 & 72 & 0 \\
\bottomrule
\end{tabular}
}
\caption{Representative routed build footprints. LUT counts are not directly
comparable across FPGA families; see the hardware evidence ledger for commands
and tool versions.}
\end{table}

\subsection{Zero-Drift Verification}
\label{sec:drift}

\begin{itemize}
\item Lucas MAC: software oracle covers 999,996 primitive operations across
  166,666 identity macros with exact bit-match return.\cite{lucas2026}
\item Lucas fast path: Tang 25K SRAM-loaded probe covers PSCALE/PCHIRAL and a
  100-period PSCALE zero-drift loop.
\item ROTC: 6-angle period closure verified on VM, RTL, and FPGA trace.
\item RPLU v2: 149-record consume profile with checksum verification
  ($\sum = \mathtt{0x0AA480E7}$).
\item RPLU2 Pad\'e: Wukong Artix-7 J11 smoke verifies five rational constant
  Pad\'e cases through the SPI-visible evaluator with no CRC error and idle
  busy status after each run.
\end{itemize}

\subsection{Power and Timing}
\label{sec:power}

The cleaned Wukong Artix-7 RPLU2PADE build routes with 72 DSP48E1 slices,
20,277 slice LUTs, 6,678 flip-flops, and no block RAM. Router convergence
occurs by iteration 9. Post-route timing reports \texttt{clk\_fast} at 36.54 MHz while
the bench image targets 2 MHz. Power and transaction-throughput measurements
remain open instrumentation tasks.

%==============================================================================
\section{Open Toolchain and OSHWA Path}
\label{sec:toolchain}
%==============================================================================

The project uses Yosys\cite{yosys} and nextpnr\cite{nextpnr} family flows
where possible:

\begin{itemize}
\item \textbf{Gowin GW5A-25A} via nextpnr-himbaechel; silicon probe target.
\item \textbf{Xilinx Artix-7} via nextpnr-xilinx/OpenXC7 and Project X-Ray;
  current Wukong integration target.
\item \textbf{Lattice ECP5} via nextpnr-ecp5 and Project Trellis; portability
  target under active bring-up.
\item \textbf{Kintex-7} via the same 7-series flow; planned full-integration
  and PCIe/DDR proof target.
\end{itemize}

Board design files (KiCad), firmware (C, Python), and tooling are released
under CERN-OHL-W, MIT, and CC0 licenses respectively. An OSHWA
self-certification package is in preparation, pending completion of the
ECP5-85K evaluation carrier board schematic.

%==============================================================================
\section{Conclusion}
\label{sec:conclusion}
%==============================================================================

The SPU-13 demonstrates that exact rational arithmetic is not a theoretical
curiosity but a practical FPGA architecture for exact geometric computation.
The present evidence base combines software oracles, RTL testbenches, and
targeted silicon proofs. With zero floating-point units, no division in the
hot datapaths, and bounded host-visible sidecar protocols, SPU-13 provides a
reproducible substrate for exact robotics, rational classification, and
algebraic approximant experiments. The open toolchain and permissive licensing
invite community contribution and independent verification.

%==============================================================================
% Bibliography
%==============================================================================

\begin{thebibliography}{9}

\bibitem{fuller1975}
R.~B. Fuller, \emph{Synergetics: Explorations in the Geometry of Thinking}.
Macmillan, 1975.

\bibitem{wildberger2005}
N.~J. Wildberger, \emph{Divine Proportions: Rational Trigonometry to Universal
Geometry}. WildEgg, 2005.

\bibitem{thomson2026}
A.~R. Thomson, ``Spread-Quadray Rotors,'' ResearchGate, May 2026.
\url{https://www.researchgate.net/profile/Andy-Ross-Thomson}

\bibitem{isa_reference}
SPU-13 ISA Reference, \texttt{knowledge/isa\_reference.md}.

\bibitem{rplu2026}
J.~Curley, ``RPLU v2: Thimble-Pad\'e Rational Approximant Pipeline over
$\mathbb{A}_{31}$,'' arXiv, 2026. \textit{(in preparation)}.

\bibitem{lucas2026}
J.~Curley, ``Zero-Drift Lucas-Prime Phinary Arithmetic for Exact Quantum
Phase Coherence on FPGA,'' arXiv, 2026. \textit{(in preparation)}.

\bibitem{yosys}
Yosys Open SYnthesis Suite. \url{https://github.com/YosysHQ/yosys}

\bibitem{nextpnr}
nextpnr: a portable FPGA place and route tool.
\url{https://github.com/YosysHQ/nextpnr}

\end{thebibliography}

\end{document}
